\documentclass[serif]{beamer}
\usepackage{etex}
\usetheme{AnnArbor}
\usecolortheme{beaver}
%\usecolortheme{dove}
\usepackage{amssymb,amsmath,amsthm}
\usepackage{graphicx}
\usepackage{tikz}
\definecolor{beetle}{RGB}{196, 30, 58}

\setbeamercolor{titlelike}{parent=structure,bg=beetle}

 \usepackage{tikz}
 \usetikzlibrary{patterns,snakes,plotmarks}
 \usepackage{multirow}
% \usetikzlibrary{shadows}
\newcommand\at[2]{\left.#1\right|_{#2}}
\usepackage{epstopdf}
\usepackage{lmodern}
\usepackage{units}
\usepackage{caption}
\usepackage{subfig}
\usepackage{qtree}
\newcommand*{\Scale}[2][4]{\scalebox{#1}{\ensuremath{#2}}}%
\DeclareCaptionLabelSeparator{horse}{:\,\,} % change according to your needs
\captionsetup{
  labelsep = horse,
  figureposition = bottom % used to get the correct vertical space between the figure and the caption
}
\begin{document}
\setbeamertemplate{caption}[numbered]
\setbeamertemplate{items}[circle]
\setbeamertemplate{enumerate items}[square]
\theoremstyle{definition}
\newtheorem*{exs}{Examples}
\newtheorem{ex}{Example}
\newtheorem*{exc}{Exercise}
%\usepackage{booktabs}
\setlength{\parindent}{0pt}
%\setbeameroption{show notes}
 \setbeamerfont{note page}{size=\tiny}
%\setbeamertemplate{note page}[plain]
%\setbeameroption{show only notes}

\title{Recursive Forecasting and Ordinal Statistical Models from Accelerometer Data}
\author{\textbf{Fatimah Alawad}}
\institute{\textbf{Ball State University}}
\date{\today}

% Print Title Page
{ % Begin Group
\usebackgroundtemplate{
\tikz\node[opacity=0.3]
{\includegraphics[height=5in,width=350pt]{unnamed.jpg}};}


\begin{frame}[plain]
\titlepage
\end{frame}
} % End Group

\begin{frame}
\frametitle{Introduction}
\begin{block}{Introduction}
\begin{itemize}
\item Since the 1980’s, accelerometers used to estimate  physical activity
\item Recently, researchers have used raw acceleration data
\item Many types of modeling used predict activity level (regression, decision trees, neural networks, etc)
\item My thesis uses random forest with recursive forecasting for improved predictive power
\item Also, I consider ordinal nature of responses
\end{itemize}
\end{block}
\end{frame}

\begin{frame}[allowframebreaks]
        \frametitle{References}
        \bibliographystyle{amsalpha}
        \bibliography{Sources.bib}
\end{frame}


 \end{document}  